% =====================================================================
\chapter{Surface de l'API et mise en œuvre}
\label{ann:api}
% =====================================================================

\section{Points d'entrée}

\begin{table}[H]
\centering
\footnotesize
\caption{Points d'entrée de l'API, avec authentification requise.}
\label{tab:api-full}
\begin{tabularx}{\textwidth}{@{}l l X@{}}
\toprule
\textbf{Méthode et chemin} & \textbf{Auth} & \textbf{Description} \\
\midrule
\api{POST /diagnostic/adresse} & oui &
  Transaction canonique, adresse $\rightarrow$ enregistrement complet \\
\api{POST /diagnostic/batch} & oui &
  Soumission d'un lot ($\leq$ 10\,000 adresses) \\
\api{GET /diagnostic/batch/\{id\}} & oui &
  État et résultats du lot \\
\api{GET /diagnostic/adresse/rapport-pdf} & oui &
  Relais vers le PDF officiel Géorisques \\
\api{GET /diagnostic/zone/building} & oui &
  Fiche BDNB au point \\
\api{GET /diagnostic/zone/buildings} & oui &
  Fiches BDNB autour du point \\
\api{GET /api/flood-alea} & oui &
  Classes TRI au point \\
\api{GET /api/hydro} & oui & Tracé hydrographique \\
\api{GET /api/hydro/extend} & oui & Extension du tracé \\
\api{GET /api/meteo} & oui & Pluie horaire et débit \\
\api{GET /api/photo} & oui & Photographie Panoramax et crédit \\
\api{GET /api/report} & oui & Rapport narratif \\
\api{GET /api/report/stream} & oui & Rapport narratif en flux \\
\api{GET /api/geocode/search} & oui & Autocomplétion d'adresse \\
\api{GET /health} & non & Sonde de vie \\
\api{GET /health/detailed} & non & Sonde détaillée \\
\bottomrule
\end{tabularx}
\end{table}

\section{Format des erreurs}

\begin{table}[H]
\centering
\small
\caption{Codes d'erreur et signification.}
\label{tab:erreurs-api}
\begin{tabularx}{\textwidth}{@{}l l X@{}}
\toprule
\textbf{Code} & \textbf{Identifiant} & \textbf{Signification} \\
\midrule
401 & — & Clé d'API absente \\
403 & — & Clé d'API invalide \\
422 & \code{adresse\_non\_trouvee} & Adresse introuvable par le géocodeur \\
422 & \code{adresse\_ambigue} & Score de géocodage inférieur à 0,4 \\
429 & \code{budget\_epuise} & Budget mensuel amont épuisé (lot uniquement) \\
502 & \code{geocodage\_indisponible} & Géocodeur injoignable ou saturé \\
502 & — & Exception inattendue (jamais de 500 HTML) \\
\bottomrule
\end{tabularx}
\end{table}

\section{Commandes d'installation}

\begin{lstlisting}[language=bash,caption={Installation et exécution en local.},
label={lst:install}]
# Backend (Python 3.11+)
cd backend
python -m venv venv && source venv/bin/activate
pip install -r requirements.txt
uvicorn app.main:app --reload --port 8000

# Frontend (Node 18+)
cd frontend
npm install
npm run dev        # http://localhost:5173
\end{lstlisting}

Le fichier d'exemple d'environnement, à copier à la racine du dépôt,
documente chaque variable nécessaire : jeton Mapbox, URL et clé Supabase, clé
Windy, clé du modèle de rédaction, identifiants Copernicus, et clés d'API
partenaires.

\section{Commandes de vérification}

\begin{lstlisting}[language=bash,caption={Exécution des suites de test.},
label={lst:tests-cmd}]
# Backend (depuis backend/)
python -m pytest tests/ -q     # 214 tests hors-ligne, < 60 s
ruff check app tests           # lint

# Frontend (depuis frontend/)
npm test                       # 143 tests Vitest
npm run build                  # tsc -b && vite build
\end{lstlisting}

\section{Vérification de l'instance déployée}

\begin{lstlisting}[language=bash,caption={Contrôles de bon fonctionnement de
l'instance en ligne.},label={lst:verif-cmd}]
BASE=https://<projet>.vercel.app
curl $BASE/health
curl "$BASE/api/flood-alea?lat=48.848&lon=2.370" | head -c 200
curl -X POST "$BASE/diagnostic/adresse" \
  -H 'Content-Type: application/json' \
  -d '{"adresse":"Quai de la Rapee, 75012 Paris"}' | head -c 200
\end{lstlisting}

Le réveil d'une fonction froide prend de 5 à 10~s, et un diagnostic complet
peut atteindre environ 15~s à froid. C'est normal sur l'offre gratuite :
appeler la sonde de vie une fois avant une démonstration suffit à préchauffer
l'instance.

\section{Reproduction des captures d'écran}

\begin{lstlisting}[language=bash,caption={Génération des captures
d'écran.},label={lst:shots}]
cd frontend && node scripts/demo-shots.mjs
# necessite le backend et le serveur vite en local
\end{lstlisting}
